\begin{proposition}[Alternate M\"obius form of prime number theorem]
\label{mu-pnt-alt}
\lean{mu_pnt_alt}
% at /home/runner/work/trough/trough/library/PrimeNumberTheoremAnd/Consequences.lean:2371.0-2429.472
\leanok
We have $\sum_{n \leq x} \mu(n)/n = o(1)$.
\end{proposition}
\begin{proof}
\uses{mu-pnt}
\leanok
As in the proof of Theorem \ref{mun}, we have
    $$ 1 = \sum_{d \leq x} \mu(d) \lfloor \frac{x}{d} \rfloor$$
    $$ = x \sum_{d \leq x} \frac{\mu(d)}{d} - \sum_{d \leq x} \mu(d) \{ \frac{x}{d} \}$$
  so it will suffice to show that
  $$ \sum_{d \leq x} \mu(d) \{ \frac{x}{d} \} = o(x).$$
  Let $N$  be a natural number.  It suffices to show that
  $$ \sum_{d \leq x} \mu(d) \{ \frac{x}{d} \} = O(x/N).$$
  if $x$ is large enough depending on $N$.
  We can split the left-hand side as the sum of
  $$ \sum_{d \leq x/N} \mu(d) \{ \frac{x}{d} \} $$
  and
  $$ \sum_{j=1}^{N-1} \sum_{x/(j+1) < d \leq x/j} \mu(d) (x/d - j).$$
  The first term is clearly $O(x/N)$.  For the second term, we can use Theorem \ref{mu-pnt}
  and summation by parts (using the fact that $x/d-j$ is monotone and bounded) to find that
  $$ \sum_{x/(j+1) < d \leq x/j} \mu(d) (x/d - j) = o(x)$$
  for any given $j$, so in particular
  $$ \sum_{x/(j+1) < d \leq x/j} \mu(d) (x/d - j) = O(x/N^2)$$
  for all $j=1,\dots,N-1$ if $x$ is large enough depending on $N$.
  Summing all the bounds, we obtain the claim.
\end{proof}
